\chapter{Reference: timers and ticks}
\label{ch:refman-timers}

\begin{aside}
Periodic time-driven messages. Use them for animation,
polling, and anything else that needs to wake up on a
schedule.
\end{aside}

\section{Forms}

\begin{itemize}[leftmargin=*]
  \item \code{Sub<Msg>::every(duration, msg)}: send \code{msg}
    every \code{duration}.
  \item \code{Cmd<Msg>::after(duration, msg)}: send
    \code{msg} once, \code{duration} from now.
\end{itemize}

\section{Working example}

\begin{lstlisting}[language=C++]
#include <maya/maya.hpp>
#include <chrono>

using namespace maya;
using namespace maya::dsl;

struct Timer {
    struct Model {
        int seconds = 0;
        bool flashed = false;
    };
    struct Tick {};
    struct Flash {};
    struct ClearFlash {};
    struct Quit {};
    using Msg = std::variant<Tick, Flash, ClearFlash, Quit>;

    static Model init() { return {}; }

    static auto update(Model m, Msg msg)
        -> std::pair<Model, Cmd<Msg>>
    {
        return std::visit(overload{
            [&](Tick) {
                m.seconds++;
                return std::pair{m, Cmd<Msg>{}};
            },
            [&](Flash) {
                m.flashed = true;
                return std::pair{m,
                    Cmd<Msg>::after(std::chrono::milliseconds(500),
                                    ClearFlash{})};
            },
            [&](ClearFlash) {
                m.flashed = false;
                return std::pair{m, Cmd<Msg>{}};
            },
            [](Quit) {
                return std::pair{Model{}, Cmd<Msg>::quit()};
            },
        }, msg);
    }

    static Element view(const Model& m) {
        auto label = std::format("elapsed: {}s", m.seconds);
        return v(
            when(m.flashed,
                text(label) | Bold | Fg<255, 220, 80>,
                text(label) | Bold),
            blank_,
            t<"f flash, q quit"> | Dim
        ) | pad<1> | border_<Round>;
    }

    static auto subscribe(const Model&) -> Sub<Msg> {
        return Sub<Msg>::batch(
            key_map<Msg>({
                {'q', Quit{}},
                {'f', Flash{}},
            }),
            Sub<Msg>::every(std::chrono::seconds(1), Tick{})
        );
    }
};

int main() { run<Timer>({.title = "timer"}); }
\end{lstlisting}

\section{Notes}

\begin{itemize}[leftmargin=*]
  \item Resolution is bounded by your \code{fps}; finer
    timers fire on the next available frame.
  \item \code{after} is "fire once" — don't use it for
    repeats.
  \item Timers do not drift; they're scheduled against a
    monotonic clock.
\end{itemize}

\section{Recap}

\begin{recap}
\begin{itemize}[leftmargin=*]
  \item \code{every} for repeats, \code{after} for one-shots.
  \item Resolution is frame-bounded.
  \item Use \code{after} to clear transient highlights.
\end{itemize}
\end{recap}

\section*{Workshop}

\begin{enumerate}[leftmargin=*]
  \item Animate a spinner using \code{every(80ms, ...)}.
  \item Show a "saved" toast for two seconds, then clear.
  \item Build a Pomodoro timer with 25-minute work blocks.
\end{enumerate}
